\section{Conclusions}

This work presents a conservative DBS formulation for PHREEQC-coupled reactive
transport with cell-based surface complexation in image-based porous media. The
formulation defines the continuous model, the discrete implementation, and the
validation sequence used to assess the solver across a range of analytical,
image-based, and multi-scale porous-media cases.

The key contribution is a cell-based representation of unresolved reactive
surface area that is consistent with the DBS micro-continuum framework. By
using the physical mobile-water fraction $\theta$ as the accumulation and
chemistry storage variable, rather than the numerical DBS porosity, the method
preserves conservative aqueous transport and well-defined local mass balances
simultaneously. The unresolved micropore surface area is
introduced as an additive intrinsic volumetric density, allowing the method to
represent reactive uptake in image phases where the internal surface area
exceeds what can be recovered from the coarse-cell boundary alone.

The reported results demonstrate that the formulation: (i) reproduces the
dilute analytical retardation limit with sub-percent bulk inventory error;
(ii) produces resolved dissolved and adsorbed fields in
transparent 2D and 3D image-based geometries; and (iii) remains computationally
tractable at application scale, with chemistry representing 36--55\% of the
total wall-clock time for the calcium surface-complexation system studied here,
with cell masking providing an effective cost control for near-solid regions.

The method closes a specific gap in existing DBS reactive transport
capabilities: conservative, unresolved surface complexation in multi-scale
image-based simulations, without requiring the reactive surface to coincide
with an explicitly resolved mesh boundary. This capability is directly relevant
to reactive porous-media problems in which selectivity and surface accessibility
control performance, including contaminant retention, radionuclide migration,
and critical-mineral recovery by sorbent-bearing porous media. The natural
extensions of this work include Strang splitting for reduced
operator-splitting error, better constraints and sensitivity studies for
intrinsic surface-area inputs and masking thresholds, and extensions to
evolving reactive area and kinetic surface chemistry.
